% Generated draft: requires mathematical and source review.
\documentclass[11pt]{article}
\usepackage[margin=25mm]{geometry}
\usepackage{fontspec}
\setmainfont{Latin Modern Roman}
\usepackage{amsmath,amssymb,mathtools,unicode-math}
\setmathfont{Latin Modern Math}
\usepackage{enumitem,needspace,xurl,hyperref}
\setlength{\parindent}{0pt}
\setlength{\parskip}{5pt}
\setlength{\emergencystretch}{3em}
\newcommand{\N}{\mathbb{N}}
\newcommand{\Z}{\mathbb{Z}}
\newcommand{\Q}{\mathbb{Q}}
\newcommand{\R}{\mathbb{R}}
\newcommand{\C}{\mathbb{C}}
\newcommand{\F}{\mathbb{F}}
\newcommand{\A}{\mathbb{A}}
\newcommand{\PP}{\mathbb{P}}
\newcommand{\E}{\mathbb{E}}
\newcommand{\eps}{\varepsilon}
\newcommand{\dd}{\,\mathrm{d}}
\begin{document}
\section*{UnsolvedMath ID: 30004347}
\textbf{Source problem code:} catalogue entry
\subsection{Definitions and mathematical statement}
For an elliptic curve \(E\) defined over a field contained in \(L\), write \(E(L)\) for its group of \(L\)-rational points. Let \(f\) be a weight-two newform equipped with an attached elliptic rigid cocycle, in the sense intended by the supplied source. Retain the source notation \(E\) for the elliptic curve involved, \(J_f^+\) and \(J_f^-\) for the two indicated cocycle components, and \(J_f^\pm[\tau]\) for their values at a real-multiplication point \(\tau\).

\textbf{Conjectural statement, subject to unresolved source specifications.} For every such attached elliptic rigid cocycle and every real-multiplication point \(\tau\) in its intended evaluation domain, both values are rational over the field denoted by \(H_\tau\):
\[
J_f^+[\tau]\in E(H_\tau)
\qquad\text{and}\qquad
J_f^-[\tau]\in E(H_\tau).
\]
The supplied record does not define \(H_\tau\), specify the ambient domain of \(\tau\), give the construction or normalization of these evaluations, or state the hypotheses relating \(f\), the cocycle, and \(E\). These are unresolved qualifications: the displayed assertion records the intended field-of-definition claim but is not a self-contained formulation until those data are supplied.
\subsection{Short English statement}
The source conjectures that both indicated values of an elliptic rigid cocycle attached to a weight-two newform, at every admissible real-multiplication point \(\tau\), lie in \(E(H_\tau)\). The evaluation domain, field \(H_\tau\), and required attachment hypotheses remain unspecified.
\subsection{Sources}
\begin{itemize}
\item Local registry: supplied UnsolvedMath source record, permanent ID \(30004347\), Oberwolfach Reports Open Problems collection, exact source code \texttt{catalogue entry}, title ``Elliptic Curve Values of Rigid Cocycles.'' The record lists proposed year 2019 and names no proposer.
\item Supplied primary report citation: \emph{Modular Forms (2020)}, Oberwolfach Reports, pp.\ 3529--3606, DOI \texttt{10.4171/owr/2019/57}. \url{https://doi.org/10.4171/owr/2019/57}. The supplied background describes these reports as documenting mathematical workshops and their open-problem discussions; no specific internal statement location is supplied.
\item Supplied literature reference: Darmon--Vonk, ``arithmetic intersections of modular geodesics.'' \url{https://doi.org/10.1016/j.jnt.2020.12.012}.
\item Supplied literature reference: Darmon--Pozzi--Vonk, ``special RM values.'' \url{https://doi.org/10.4171/JEMS/1344}.
\item Supplied literature assessment, checked \textbf{2026-08-22}: the general field-of-definition conjecture for real-multiplication values of the \(f\)-components of rigid meromorphic cocycles is labelled open, with current theorems described as covering special cocycles such as the Dedekind--Rademacher class. This is supplied dated curation derived from the Oberwolfach report collection, not an independently verified literature assessment; its status and cited claims require independent verification. The report above is also included among the supplied literature sources.
\end{itemize}
\end{document}
